\documentclass[12pt]{article}

% =========================
% Geometry & core packages
% =========================
\usepackage[a4paper,margin=0.8in]{geometry}
\usepackage{amsmath,amssymb,amsthm,mathtools}
\usepackage{bm}
\usepackage{array,booktabs}
\usepackage{enumitem}
\usepackage{microtype}
\usepackage{xcolor}
\usepackage{hyperref}
\usepackage{fancyhdr}
\usepackage{draftwatermark}
\usepackage{titlesec}

% =========================
% Watermark
% =========================
\SetWatermarkText{Unofficial — QRS@NTU — qrsntu.org}
\SetWatermarkScale{1}
\SetWatermarkLightness{0.99}

% =========================
% Metadata
% =========================
\newcommand{\CourseCode}{MH1101}
\newcommand{\CourseName}{Calculus II}
\newcommand{\DocType}{Past Year Paper Solutions}
\newcommand{\ExamMeta}{AY 2018/2019, Semester 2}
\newcommand{\ExamMetaShort}{Midterm AY18-19 Sem 2}

\title{
\Huge \CourseCode\ \CourseName \\[0.3em]
\Large \DocType \\[0.3em]
\large \ExamMeta
}
\author{
  \textit{\href{https://qrsntu.org}{Quantitative Research Society @NTU}}
}
\date{\today}

% =========================
% Header/footer
% =========================
\fancypagestyle{main}{
  \fancyhf{}
  \fancyhead[L]{\CourseCode\ \CourseName\ --\ \ExamMetaShort}
  \fancyhead[R]{\href{https://qrsntu.org}{QRS@NTU}}
  \fancyfoot[C]{\thepage}
  \renewcommand{\headrulewidth}{0.4pt}
  \renewcommand{\footrulewidth}{0pt}
}

\fancypagestyle{plainfirst}{
  \fancyhf{}
  \renewcommand{\headrulewidth}{0pt}
  \renewcommand{\footrulewidth}{0pt}
  \fancyfoot[C]{\scriptsize\textcolor{gray}{\textit{For academic reference only. This document is provided for educational purposes and does not constitute an official question paper, answer key, or any formally authorised academic material. Its content is not guaranteed to be complete or accurate.}}}
}

% =========================
% Convenience macros
% =========================
\newcommand{\dd}{\,\mathrm{d}}
\newcommand{\ee}{\mathrm{e}}
\newcommand{\R}{\mathbb{R}}
\newcommand{\N}{\mathbb{N}}

% =========================
% Overview block
% =========================
\newenvironment{overview}{
  \par\bigskip
  \begin{center}
    \rule{0.92\linewidth}{0.6pt}\\[-0.2em]
    \textbf{Overview}\\[-0.2em]
    \rule{0.92\linewidth}{0.6pt}
  \end{center}
  \vspace{-0.3em}
  \begin{itemize}[leftmargin=1.6em, itemsep=0.2em, topsep=0.2em]
}{
  \end{itemize}
  \par\bigskip
}

\titleformat{\subsubsection}[block]
{\normalfont\large\bfseries}
{}
{0pt}
{}

\titleformat{\paragraph}[block]
{\normalfont\normalsize\bfseries}
{}
{0pt}
{}

\titlespacing*{\paragraph}
{0pt}
{1.2ex}
{0.6ex}

\setlist[itemize]{itemsep=0.2em, topsep=0.2em}

\setcounter{secnumdepth}{-1}
\setcounter{tocdepth}{4}

\begin{document}

\maketitle
\thispagestyle{plainfirst}
\pagestyle{main}

\begin{overview}
  \item Fundamental Theorem of Calculus and differentiation of integrals with variable limits.
  \item Riemann sums and their interpretation as definite integrals.
  \item Applications of integration: volume of solids of revolution.
  \item Gaussian-type improper integrals and scaling substitutions.
  \item Techniques of integration: substitution, trigonometric identities, and integration by parts.
  \item Recursive evaluation of parameter-dependent definite integrals.
\end{overview}

\newpage
\tableofcontents
\newpage

\section{Question 1 \hfill (16 marks)}


\subsection{(a) Differentiate an Integral with Variable Limits \hfill (8 marks)}

Calculate the derivative
\[
\frac{\dd}{\dd x}\int_{4x}^{2x}\cos\!\left(\ee^{2t}\right)\dd t.
\]

\subsubsection{Solution}

\paragraph{Method 1: Leibniz rule directly}
Let
\[
F(x)=\int_{4x}^{2x}\cos\!\left(\ee^{2t}\right)\dd t.
\]
Using
\[
\frac{\dd}{\dd x}\int_{u(x)}^{v(x)} f(t)\,\dd t
=f(v(x))v'(x)-f(u(x))u'(x),
\]
with
\[
f(t)=\cos\!\left(\ee^{2t}\right),\qquad u(x)=4x,\qquad v(x)=2x,
\]
we obtain
\[
F'(x)=\cos\!\left(\ee^{2(2x)}\right)\cdot 2-\cos\!\left(\ee^{2(4x)}\right)\cdot 4.
\]
Hence
\[
F'(x)=2\cos\!\left(\ee^{4x}\right)-4\cos\!\left(\ee^{8x}\right).
\]
Therefore
\[
\boxed{\frac{\dd}{\dd x}\int_{4x}^{2x}\cos\!\left(\ee^{2t}\right)\dd t
=2\cos\!\left(\ee^{4x}\right)-4\cos\!\left(\ee^{8x}\right).}
\]

\paragraph{Method 2: Rewrite using an auxiliary accumulation function}
Define
\[
G(y)=\int_0^y \cos\!\left(\ee^{2t}\right)\dd t.
\]
Then
\[
\int_{4x}^{2x}\cos\!\left(\ee^{2t}\right)\dd t
=G(2x)-G(4x).
\]
By the Fundamental Theorem of Calculus,
\[
G'(y)=\cos\!\left(\ee^{2y}\right).
\]
Thus
\[
\frac{\dd}{\dd x}\bigl(G(2x)-G(4x)\bigr)
=2G'(2x)-4G'(4x)
\boxed{=2\cos\!\left(\ee^{4x}\right)-4\cos\!\left(\ee^{8x}\right).}
\]


\subsubsection{Marking scheme}
\begin{itemize}
\item 2 marks: Correct use of the Leibniz rule or an equivalent FTC decomposition.
\item 2 marks: Correct evaluation of the integrand at \(t=2x\) and \(t=4x\).
\item 2 marks: Correct differentiation of the upper and lower limits.
\item 2 marks: Correct final expression \(2\cos(\ee^{4x})-4\cos(\ee^{8x})\).
\end{itemize}

\subsection{(b) Write the Limit as a Definite Integral \hfill (8 marks)}

Write the following limit as a definite integral (do not evaluate it):
\[
\lim_{n\to\infty}\sum_{i=1}^{n}\sqrt{\frac{ni+3i^2}{n^4}}.
\]

\subsubsection{Solution}

\paragraph{Method 1: Recognise the Riemann sum on \([0,1]\)}
First simplify the summand:
\[
\sqrt{\frac{ni+3i^2}{n^4}}
=\sqrt{\frac{ni}{n^4}+\frac{3i^2}{n^4}}
=\sqrt{\frac{i}{n^3}+3\frac{i^2}{n^4}}.
\]
Now factor out \(1/n\):
\[
\sqrt{\frac{i}{n^3}+3\frac{i^2}{n^4}}
=\frac{1}{n}\sqrt{\frac{i}{n}+3\left(\frac{i}{n}\right)^2}.
\]
Hence
\[
\sum_{i=1}^{n}\sqrt{\frac{ni+3i^2}{n^4}}
=\sum_{i=1}^{n}\sqrt{\frac{i}{n}+3\left(\frac{i}{n}\right)^2}\,\frac{1}{n}.
\]
This is a right-endpoint Riemann sum for
\[
f(x)=\sqrt{x+3x^2}
\]
on the interval \([0,1]\), with
\[
x_i=\frac{i}{n},\qquad \Delta x=\frac{1}{n}.
\]
Therefore
\[
\lim_{n\to\infty}\sum_{i=1}^{n}\sqrt{\frac{ni+3i^2}{n^4}}
\quad \boxed{=\int_0^1 \sqrt{x+3x^2}\,\dd x}\,.
\]

\newpage
\paragraph{Method 2: Introduce \(x_i=i/n\) from the start}
Set
\[
x_i=\frac{i}{n}.
\]
Then
\[
i=nx_i,
\qquad
ni+3i^2=n(nx_i)+3(nx_i)^2=n^2x_i+3n^2x_i^2=n^2(x_i+3x_i^2).
\]
So
\[
\sqrt{\frac{ni+3i^2}{n^4}}
=\sqrt{\frac{n^2(x_i+3x_i^2)}{n^4}}
=\frac{1}{n}\sqrt{x_i+3x_i^2}.
\]
Hence
\[
\sum_{i=1}^{n}\sqrt{\frac{ni+3i^2}{n^4}}
=\sum_{i=1}^{n}\sqrt{x_i+3x_i^2}\,\frac{1}{n}.
\]
Therefore the limit is
\[
\boxed{\int_0^1 \sqrt{x+3x^2}\,\dd x.}
\]

\subsubsection{Marking scheme}
\begin{itemize}
\item 2 marks: Correct algebraic rewriting of the summand into the form \(f(i/n)\cdot (1/n)\).
\item 2 marks: Correct identification of the interval \([0,1]\).
\item 2 marks: Correct identification of the integrand \(\sqrt{x+3x^2}\).
\item 2 marks: Correct final definite integral, without evaluation.
\end{itemize}

\newpage

\section{Question 2 \hfill (16 marks)}

\subsection{(a) Volume of Revolution about \(x=4\) \hfill (10 marks)}

Let \(R\) be the region bounded by the curves \(y=x^2\) and \(y=4x-x^2\). Compute the volume of the solid formed by revolving \(R\) about the line \(x=4\).

\subsubsection{Solution}

\paragraph{Method 1: Cylindrical shells}
First find the points of intersection:
\[
x^2=4x-x^2
\quad\Longrightarrow\quad
2x^2-4x=0
\quad\Longrightarrow\quad
2x(x-2)=0.
\]
Thus
\[
x=0,\qquad x=2.
\]
For \(0\le x\le 2\), the upper curve is
\[
y=4x-x^2
\]
and the lower curve is
\[
y=x^2.
\]
Since the axis of rotation is the vertical line \(x=4\), the shell method is natural.

A shell at position \(x\) has radius
\[
r(x)=4-x
\]
and height
\[
h(x)=(4x-x^2)-x^2=4x-2x^2.
\]
Hence
\[
V=2\pi\int_0^2 (4-x)(4x-2x^2)\,\dd x.
\]
Expand:
\[
(4-x)(4x-2x^2)=16x-12x^2+2x^3.
\]
So
\begin{align*}
V
&=2\pi\int_0^2 (16x-12x^2+2x^3)\,\dd x \\
&=2\pi\left[8x^2-4x^3+\frac{x^4}{2}\right]_0^2 \\
&=2\pi(32-32+8) \\
&=16\pi.
\end{align*}
Therefore
\[
\boxed{V=16\pi.}
\]

\paragraph{Method 2: Washers with respect to \(y\)}
Rewrite the curves in terms of \(x\) as functions of \(y\).

From
\[
y=x^2,
\]
we obtain
\[
x=\sqrt{y}
\]
in the relevant region.

From
\[
y=4x-x^2,
\]
we have
\[
x^2-4x+y=0,
\]
so
\[
x=2\pm \sqrt{4-y}.
\]
The region extends from \(y=0\) to \(y=4\). A horizontal slice gives a washer about \(x=4\).

For \(0\le y\le 1\), the slice runs from
\[
x=2-\sqrt{4-y}
\quad\text{to}\quad
x=\sqrt{y}.
\]
Thus
\[
R(y)=4-\bigl(2-\sqrt{4-y}\bigr)=2+\sqrt{4-y},
\qquad
r(y)=4-\sqrt{y}.
\]
For \(1\le y\le 4\), the slice runs from
\[
x=2-\sqrt{4-y}
\quad\text{to}\quad
x=2+\sqrt{4-y}.
\]
Thus
\[
R(y)=2+\sqrt{4-y},
\qquad
r(y)=2-\sqrt{4-y}.
\]
Therefore
\begin{align*}
V
&=\pi\int_0^1\Bigl[(2+\sqrt{4-y})^2-(4-\sqrt{y})^2\Bigr]\dd y \\
&\quad +\pi\int_1^4\Bigl[(2+\sqrt{4-y})^2-(2-\sqrt{4-y})^2\Bigr]\dd y.
\end{align*}
Although this setup is correct, it is substantially more cumbersome to evaluate than the shell method. Evaluating the two definite integrals yields the same result:
\[
\boxed{V=16\pi.}
\]

\subsubsection{Marking scheme}
\begin{itemize}
\item 2 marks: Correct intersection points \(x=0\) and \(x=2\).
\item 2 marks: Correct identification of a valid volume method.
\item 2 marks: Correct radius and height, or correct washer radii.
\item 2 marks: Correct integral setup with correct limits.
\item 2 marks: Correct evaluation to \(V=16\pi\).
\end{itemize}

\subsection{(b) Scaled Gaussian Integral \hfill (6 marks)}

Given that
\[
\int_{-\infty}^{\infty}\ee^{-x^2}\dd x=\sqrt{\pi},
\]
evaluate
\[
\int_{-\infty}^{\infty}\ee^{-kx^2}\dd x
\qquad \text{for } k>0.
\]
Justify your answer.

\subsubsection{Solution}
\paragraph{Method 1: Direct substitution}
Let
\[
I=\int_{-\infty}^{\infty}\ee^{-kx^2}\dd x,
\qquad k>0.
\]
Use the substitution
\[
u=\sqrt{k}\,x.
\]
Then
\[
\dd u=\sqrt{k}\,\dd x,
\qquad
\dd x=\frac{1}{\sqrt{k}}\dd u.
\]
Since \(k>0\),
\[
x\to -\infty \implies u\to -\infty,
\qquad
x\to \infty \implies u\to \infty.
\]
Therefore
\begin{align*}
I
&=\int_{-\infty}^{\infty}\ee^{-u^2}\frac{1}{\sqrt{k}}\dd u \\
&=\frac{1}{\sqrt{k}}\int_{-\infty}^{\infty}\ee^{-u^2}\dd u \\
&=\frac{1}{\sqrt{k}}\sqrt{\pi}.
\end{align*}
Hence
\[
\boxed{\int_{-\infty}^{\infty}\ee^{-kx^2}\dd x=\sqrt{\frac{\pi}{k}}.}
\]

\newpage
\paragraph{Method 2: Split at \(0\) before the change of variable}
Write
\[
\int_{-\infty}^{\infty}\ee^{-kx^2}\dd x
=
\int_{-\infty}^{0}\ee^{-kx^2}\dd x
+
\int_{0}^{\infty}\ee^{-kx^2}\dd x.
\]
Now use
\[
u=\sqrt{k}\,x
\]
in each integral. Then
\[
\dd x=\frac{1}{\sqrt{k}}\dd u.
\]
So
\begin{align*}
\int_{-\infty}^{0}\ee^{-kx^2}\dd x
&=\frac{1}{\sqrt{k}}\int_{-\infty}^{0}\ee^{-u^2}\dd u,\\
\int_{0}^{\infty}\ee^{-kx^2}\dd x
&=\frac{1}{\sqrt{k}}\int_{0}^{\infty}\ee^{-u^2}\dd u.
\end{align*}
Adding,
\begin{align*}
\int_{-\infty}^{\infty}\ee^{-kx^2}\dd x
&=\frac{1}{\sqrt{k}}\left(\int_{-\infty}^{0}\ee^{-u^2}\dd u+\int_{0}^{\infty}\ee^{-u^2}\dd u\right) \\
&=\frac{1}{\sqrt{k}}\int_{-\infty}^{\infty}\ee^{-u^2}\dd u \\
&=\frac{\sqrt{\pi}}{\sqrt{k}}.
\end{align*}
Thus
\[
\boxed{\int_{-\infty}^{\infty}\ee^{-kx^2}\dd x=\sqrt{\frac{\pi}{k}}.}
\]

\subsubsection{Marking scheme}
\begin{itemize}
\item 2 marks: Correct substitution \(u=\sqrt{k}\,x\).
\item 2 marks: Correct transformation of the differential and limits.
\item 2 marks: Correct final value \(\sqrt{\pi/k}\), with justification.
\end{itemize}

\newpage

\section{Question 3 \hfill (18 marks)}

\subsection{(a) Trigonometric Integral \hfill (6 marks)}

Evaluate
\[
\int \sin^5(3x)\cos^4(3x)\,\dd x.
\]

\subsubsection{Solution}
\paragraph{Method 1: Save one odd sine factor}
Since the power of sine is odd, write
\[
\sin^5(3x)=\sin^4(3x)\sin(3x)=\bigl(1-\cos^2(3x)\bigr)^2\sin(3x).
\]
Hence
\[
\int \sin^5(3x)\cos^4(3x)\,\dd x
=
\int \bigl(1-\cos^2(3x)\bigr)^2\cos^4(3x)\sin(3x)\,\dd x.
\]
Let
\[
u=\cos(3x).
\]
Then
\[
\dd u=-3\sin(3x)\,\dd x,
\qquad
\sin(3x)\,\dd x=-\frac{1}{3}\dd u.
\]
Therefore
\begin{align*}
\int \sin^5(3x)\cos^4(3x)\,\dd x
&=-\frac{1}{3}\int (1-u^2)^2u^4\,\dd u \\
&=-\frac{1}{3}\int (u^4-2u^6+u^8)\,\dd u \\
&=-\frac{1}{3}\left(\frac{u^5}{5}-\frac{2u^7}{7}+\frac{u^9}{9}\right)+C.
\end{align*}
Substituting back,
\[
\boxed{\int \sin^5(3x)\cos^4(3x)\,\dd x
=-\frac{\cos^5(3x)}{15}
+\frac{2\cos^7(3x)}{21}
-\frac{\cos^9(3x)}{27}+C.}
\]

\newpage
\paragraph{Method 2: Use \(u=\sin(3x)\)}
Write
\[
\cos^4(3x)=\cos^3(3x)\cos(3x)=\cos^2(3x)\cos(3x)\cos(3x).
\]
Using
\[
\cos^2(3x)=1-\sin^2(3x),
\]
we obtain
\[
\sin^5(3x)\cos^4(3x)
=\sin^5(3x)\bigl(1-\sin^2(3x)\bigr)^2\cos(3x).
\]
Hence
\[
\int \sin^5(3x)\cos^4(3x)\,\dd x
=
\int \sin^5(3x)\bigl(1-\sin^2(3x)\bigr)^2\cos(3x)\,\dd x.
\]
Let
\[
u=\sin(3x).
\]
Then
\[
\dd u=3\cos(3x)\,\dd x,
\qquad
\cos(3x)\,\dd x=\frac{1}{3}\dd u.
\]
So
\begin{align*}
\int \sin^5(3x)\cos^4(3x)\,\dd x
&=\frac{1}{3}\int u^5(1-u^2)^2\,\dd u \\
&=\frac{1}{3}\int (u^5-2u^7+u^9)\,\dd u \\
&=\frac{u^6}{18}-\frac{u^8}{12}+\frac{u^{10}}{30}+C.
\end{align*}
Thus an equivalent antiderivative is
\[
\boxed{\int \sin^5(3x)\cos^4(3x)\,\dd x
=\frac{\sin^6(3x)}{18}-\frac{\sin^8(3x)}{12}+\frac{\sin^{10}(3x)}{30}+C.}
\]

\subsubsection{Marking scheme}
\begin{itemize}
\item 2 marks: Correct trigonometric rewriting to prepare for substitution.
\item 2 marks: Correct substitution and polynomial integration.
\item 2 marks: Correct final antiderivative.
\end{itemize}

\newpage
\subsection{(b) Logarithmic Integral \hfill (6 marks)}

Evaluate
\[
\int \frac{(\ln x)^2}{x^2}\,\dd x.
\]

\subsubsection{Solution}
\paragraph{Method 1: Integration by parts twice}
Let
\[
I=\int \frac{(\ln x)^2}{x^2}\,\dd x.
\]
Take
\[
u=(\ln x)^2,\qquad \dd v=x^{-2}\,\dd x.
\]
Then
\[
\dd u=2(\ln x)x^{-1}\,\dd x,
\qquad
v=-x^{-1}=-\frac{1}{x}.
\]
Hence
\begin{align*}
I
&=uv-\int v\,\dd u \\
&=-\frac{(\ln x)^2}{x}+2\int \frac{\ln x}{x^2}\,\dd x.
\end{align*}
Now compute
\[
J=\int \frac{\ln x}{x^2}\,\dd x
\]
by integration by parts again. Let
\[
U=\ln x,\qquad \dd V=x^{-2}\,\dd x.
\]
Then
\[
\dd U=\frac{1}{x}\,\dd x,
\qquad
V=-\frac{1}{x}.
\]
Therefore
\begin{align*}
J
&=UV-\int V\,\dd U \\
&=-\frac{\ln x}{x}+\int x^{-2}\,\dd x \\
&=-\frac{\ln x}{x}-\frac{1}{x}+C.
\end{align*}
Substituting into \(I\),
\begin{align*}
I
&=-\frac{(\ln x)^2}{x}+2\left(-\frac{\ln x}{x}-\frac{1}{x}\right)+C \\
&=-\frac{(\ln x)^2+2\ln x+2}{x}+C.
\end{align*}
Thus
\[
\boxed{\int \frac{(\ln x)^2}{x^2}\,\dd x
=-\frac{(\ln x)^2+2\ln x+2}{x}+C.}
\]

\newpage
\paragraph{Method 2: Differentiate a structured guess}
Seek an antiderivative of the form
\[
F(x)=-\frac{A(\ln x)}{x},
\]
where \(A\) is a polynomial. Try
\[
F(x)=-\frac{(\ln x)^2+2\ln x+2}{x}.
\]
Differentiate:
\begin{align*}
F'(x)
&=-\left(\frac{2\ln x+2}{x}\cdot \frac{1}{x}\right)
+\frac{(\ln x)^2+2\ln x+2}{x^2} \\
&=\frac{-(2\ln x+2)+(\ln x)^2+2\ln x+2}{x^2} \\
&=\frac{(\ln x)^2}{x^2}.
\end{align*}
Hence
\[
\boxed{\int \frac{(\ln x)^2}{x^2}\,\dd x
=-\frac{(\ln x)^2+2\ln x+2}{x}+C.}
\]

\subsubsection{Marking scheme}
\begin{itemize}
\item 2 marks: Correct integration-by-parts setup or equivalent valid method.
\item 2 marks: Correct intermediate simplification.
\item 2 marks: Correct final antiderivative.
\end{itemize}

\newpage
\subsection{(c) Parameter-Dependent Definite Integral \hfill (6 marks)}

Evaluate
\[
\int_0^1 x^a(1-x)^b\,\dd x,
\qquad \text{where \(a\) and \(b\) are non-negative integers.}
\]
Express your answer in terms of \(a\) and \(b\).

\subsubsection*{Solution}

Let
\[
I(a,b)=\int_0^1 x^a(1-x)^b\,\dd x,
\qquad a,b\in \{0,1,2,\dots\}.
\]

\paragraph{Method 1: Recursive integration by parts}

Assume first that \(a>0\). Take
\[
u=x^a,\qquad \dd v=(1-x)^b\,\dd x.
\]
Then
\[
\dd u=ax^{a-1}\,\dd x,
\qquad
v=-\frac{(1-x)^{b+1}}{b+1}.
\]
Hence
\begin{align*}
I(a,b)
&=\left[-\frac{x^a(1-x)^{b+1}}{b+1}\right]_0^1
+\frac{a}{b+1}\int_0^1 x^{a-1}(1-x)^{b+1}\,\dd x.
\end{align*}
The boundary term is zero, so
\[
I(a,b)=\frac{a}{b+1}I(a-1,b+1).
\]

Repeating this reduction \(a\) times,
\[
I(a,b)
=\frac{a!}{(b+1)(b+2)\cdots(b+a)}I(0,b+a).
\]

Now
\[
I(0,b+a)=\int_0^1(1-x)^{a+b}\,\dd x
=\frac{1}{a+b+1}.
\]

Also,
\[
(b+1)(b+2)\cdots(b+a)=\frac{(a+b)!}{b!}.
\]

Therefore
\[
I(a,b)=\frac{a!}{(a+b)!/b!}\cdot\frac{1}{a+b+1}
=\frac{a!\,b!}{(a+b+1)!}.
\]

Thus
\[
\boxed{\int_0^1 x^a(1-x)^b\,\dd x=\frac{a!\,b!}{(a+b+1)!}.}
\]

\paragraph{Method 2: Binomial expansion to obtain a finite-sum form}

Using the binomial theorem,
\[
(1-x)^b=\sum_{k=0}^{b}\binom{b}{k}(-1)^k x^k.
\]
Hence
\begin{align*}
I(a,b)
&=\int_0^1 x^a(1-x)^b\,\dd x \\
&=\int_0^1 x^a\sum_{k=0}^{b}\binom{b}{k}(-1)^k x^k\,\dd x \\
&=\sum_{k=0}^{b}\binom{b}{k}(-1)^k\int_0^1 x^{a+k}\,\dd x \\
&=\sum_{k=0}^{b}\binom{b}{k}\frac{(-1)^k}{a+k+1}.
\end{align*}

Thus the integral may be written equivalently as
\[
\boxed{
\int_0^1 x^a(1-x)^b\,\dd x
=\sum_{k=0}^{b}\binom{b}{k}\frac{(-1)^k}{a+k+1}.
}
\]


\paragraph{Method 3: Coefficient comparison}

Let \(n=a+b\). Consider
\[
(tx+1-x)^n.
\]

Expanding,
\[
(tx+1-x)^n
=\sum_{k=0}^{n}\binom{n}{k}t^k x^k(1-x)^{n-k}.
\]

Integrating from \(0\) to \(1\),
\[
\sum_{k=0}^{n}\binom{n}{k}t^k
\int_0^1 x^k(1-x)^{n-k}\,\dd x
=\int_0^1 (tx+1-x)^n\,\dd x.
\]

Now
\[
(tx+1-x)^n=(1-(1-t)x)^n.
\]

Hence
\begin{align*}
\int_0^1 (tx+1-x)^n\,\dd x
&=\int_0^1 (1-(1-t)x)^n\,\dd x \\
&=\frac{1-t^{n+1}}{(n+1)(1-t)}
=\frac{1+t+\cdots+t^n}{n+1}.
\end{align*}

Comparing coefficients of \(t^a\),
\[
\binom{n}{a}\int_0^1 x^a(1-x)^{n-a}\,\dd x
=\frac{1}{n+1}.
\]

Thus
\[
\int_0^1 x^a(1-x)^b\,\dd x
=\frac{1}{(a+b+1)\binom{a+b}{a}}
=\frac{a!\,b!}{(a+b+1)!}.
\]

\paragraph{Method 4: Two-variable substitution}

Consider
\[
A=\int_0^\infty\int_0^\infty e^{-(u+v)}u^a v^b\,\dd u\,\dd v.
\]

Since the integrals factor,
\[
A=\left(\int_0^\infty e^{-u}u^a\,\dd u\right)
\left(\int_0^\infty e^{-v}v^b\,\dd v\right)
=a!\,b!.
\]

Now introduce the change of variables
\[
u=xt,\qquad v=(1-x)t,
\qquad t=u+v.
\]

The Jacobian of this transformation is \(t\), so
\[
e^{-(u+v)}u^a v^b\,\dd u\,\dd v
=e^{-t}t^{a+b+1}x^a(1-x)^b\,\dd x\,\dd t.
\]

Hence
\[
A=\int_0^\infty e^{-t}t^{a+b+1}\,\dd t
\int_0^1 x^a(1-x)^b\,\dd x.
\]

But
\[
\int_0^\infty e^{-t}t^{a+b+1}\,\dd t=(a+b+1)!.
\]

Therefore
\[
I(a,b)=\frac{a!\,b!}{(a+b+1)!}.
\]

\paragraph{Method 5: Symmetry and simplex interpretation}

Consider the triangular region
\[
T=\{(u,v)\mid u\ge0,\ v\ge0,\ u+v\le1\}.
\]

Then
\begin{align*}
\int_T u^a v^b\,\dd u\,\dd v
&=\int_0^1\int_0^{1-u}u^a v^b\,\dd v\,\dd u \\
&=\frac{1}{b+1}\int_0^1 u^a(1-u)^{b+1}\,\dd u.
\end{align*}

Thus
\[
\int_0^1 x^a(1-x)^b\,\dd x
=(b+1)\int_T u^a v^b\,\dd u\,\dd v.
\]

Evaluating the simplex integral recursively yields
\[
\int_T u^a v^b\,\dd u\,\dd v
=\frac{a!\,b!}{(a+b+2)!}.
\]

Therefore
\[
\int_0^1 x^a(1-x)^b\,\dd x
=\frac{a!\,b!}{(a+b+1)!}.
\]

\bigskip

Hence, by any of the above methods,
\[
\boxed{\int_0^1 x^a(1-x)^b\,\dd x=\frac{a!\,b!}{(a+b+1)!}.}\]

\subsubsection{Marking scheme}
\begin{itemize}
\item 2 marks: Correct recurrence setup by integration by parts, or equivalent valid expansion.
\item 2 marks: Correct reduction or evaluation process.
\item 2 marks: Correct final formula \(\dfrac{a!\,b!}{(a+b+1)!}\).
\end{itemize}

\end{document}